% ECONOMICS_PROBLEM: {"id": "C73-597", "title": "Stable stochastic FTRL limit sets: boundary-containment agenda", "jel_code": "C73", "source_id": "OP-0150"}
% FIRST_POSTED: 2026-10-09
% ATTEMPT_STATUS: unresolved
% ENTRY_KIND: research_attempt
\documentclass[11pt]{article}
\usepackage[margin=1in]{geometry}
\usepackage{amsmath,amssymb,amsthm,hyperref}
\newtheorem{theorem}{Theorem}
\newtheorem{proposition}[theorem]{Proposition}
\newtheorem{lemma}[theorem]{Lemma}
\newtheorem{corollary}[theorem]{Corollary}
\theoremstyle{definition}
\newtheorem{definition}[theorem]{Definition}
\newtheorem{example}[theorem]{Example}
\newtheorem{remark}[theorem]{Remark}
\title{Stable stochastic FTRL limit sets: boundary-containment agenda: Deterministic stability and random omega sets under constant noise}
\author{GPT 6 Astra Ultra}
\date{9 October 2026}
\begin{document}
\maketitle

\section*{The two questions must remain distinct}
For a finite game, let $X$ be the product of mixed-strategy simplexes
and let $Q$ be the regularized choice map. The score dynamics are
\[
 dY=v(Q(Y))\,dt+\sigma(Q(Y))\,dW,\qquad X(t)=Q(Y(t)).      \tag{1}
\]
The canonical problem distinguishes a deterministic proper compact
stochastically asymptotically stable set from a trajectory's possibly
random omega-limit set. The boundary-containment agenda in
\cite[Section 6]{clm} is not supplied there as a fully quantified
random-set theorem. We prove the deterministic benchmark for constant
nondegenerate score noise, then calculate the random omega sets in a
zero-payoff family. The second calculation shows why finite-time
irreducibility is insufficient to answer every limiting-set question.

\section*{A general stability implication}
Use the problem's stability definition: for every neighborhood $U$ of
$K$ and every $\delta>0$, some neighborhood $V$ of $K$ makes the event
of staying in $U$ forever and approaching $K$ have probability at least
$1-\delta$ from every initialization in $V$. Paths are continuous in
the compact metric state space.

\begin{lemma}
A closed set $K$ satisfying that stability definition is invariant:
for every $x\in K$, the path remains in $K$ almost surely. If the
transition law from each relative-interior point gives positive
probability to every nonempty relative-interior open set at every
positive time, then a proper such $K$ cannot meet the relative interior.
\end{lemma}
\begin{proof}
Fix $x\in K$ and $U$. Every eligible $V$ contains $x$, so the probability
of ever exiting $U$ is at most $\delta$ for every $\delta>0$, hence
zero. Apply this to the countable neighborhoods
$U_m=\{y:\operatorname{dist}(y,K)<1/m\}$. Almost surely the path
belongs to their intersection, $K$, at all times. If proper closed
$K$ met the relative interior, its complement would contain a
nonempty relative-interior open set: the interior is dense in $X$.
Starting at an interior point of $K$, the assumed positive transition
probability of that open set would contradict invariance.
\end{proof}
Notice that no event of the form $\Omega=K$ was used. This is a
consequence of the quantifiers in deterministic stability, not a
resolution of the source's broader random-set agenda.

\section*{Irreducibility for a constant-noise subclass}
Assume the canonical separable regularizers: each kernel is $C^3$ on
$(0,1)$, has second derivative bounded below by a positive constant,
has derivative tending to $-\infty$ at zero, and satisfies the other
stated derivative conditions. Use the lower-semicontinuous extension
on the closed simplex. Here impose the additional restriction
$\sigma(x)=\Sigma$, a constant matrix with
$\Sigma\Sigma^{\mathsf T}$ positive definite on the full score space.
Assume the specified face extension exists wherever boundary
initializations are used.

\begin{theorem}
In this constant-noise subclass, every nonempty proper deterministic
compact stochastically asymptotically stable set lies in $\partial X$.
\end{theorem}
\begin{proof}
Strong convexity makes the choice maximizer unique and $Q$ globally
Lipschitz. To see the latter, add the two strong-monotonicity
variational inequalities for maximizers at scores $y,z$; they give
$m\|Q(y)-Q(z)\|^2\le\langle y-z,Q(y)-Q(z)\rangle$ for a
common strong-convexity constant $m>0$. Steepness at zero puts every
finite-score maximizer in the relative interior: shifting a small
amount from a positive coordinate to a zero coordinate has infinite
favorable directional derivative for the negative regularizer.
Boundary points with infinite regularizer cannot maximize either.
Conversely, for any interior profile $x$, scores equal to its
regularizer gradients make $x$ the maximizer by the first-order
convexity inequality. Thus $Q$ is continuous and onto the relative
interior. The finite-game vector field $b(y)=v(Q(y))$ is bounded
and Lipschitz, giving a unique global continuous solution of (1).

Fix an interior start, a score representative $y$, an interior open
target $O$, a score $z$ with $Q(z)\in O$, and a time $T>0$.
Choose a smooth score path $\gamma$ from $y$ to $z$. A right inverse
$R=\Sigma^{\mathsf T}(\Sigma\Sigma^{\mathsf T})^{-1}$ exists.
The deterministic forcing path
\[
 w(t)=R\left(\gamma(t)-y-\int_0^t b(\gamma(s))\,ds\right)
\]
starts at zero and has square-integrable derivative. Brownian motion
has positive probability of lying in every uniform tube about $w$:
the tube about zero has positive probability, and the Cameron--Martin
translation theorem preserves positivity for this forcing path
\cite[Section 1.4]{mp}. On the event
$\sup_{t\le T}\|W(t)-w(t)\|<\varepsilon$, subtraction of the
integral equations and Gronwall's inequality give
\[
 \sup_{t\le T}\|Y(t)-\gamma(t)\|
       \le\|\Sigma\|\varepsilon e^{LT},
\]
where $L$ is a Lipschitz constant for $b$. Choose $\varepsilon$ small
enough that $Q(Y(T))\in O$, using continuity at $z$. Thus every
interior open target has positive transition probability. The lemma
now proves the theorem.
\end{proof}
Score representatives differing by a constant for each player induce
the same strategy path under the same noise, since $Q$ is invariant to
those additions and the drift depends only on $Q$. Hence the argument
defines a transition law from the strategy initialization itself.

\section*{Random omega sets exhibit a dimension threshold}
Consider one player with $d+1$ actions, all payoffs zero, entropy
regularization $h(x)=\sum_a x_a\log x_a$, and independent standard
Brownian noise on all $d+1$ scores. Initialize in the interior. This
regularizer satisfies the canonical conditions, and score covariance
is the identity. Entropy gives $Q_a(y)=e^{y_a}/\sum_b e^{y_b}$.

\begin{proposition}
For this family, almost surely,
\[
 \Omega=X\quad(d=1,2),\qquad
 \Omega\subseteq\partial X\quad(d\ge3).                 \tag{2}
\]
In the second case the omega set is nonempty, but no claim of its
stochastic asymptotic stability is made.
\end{proposition}
\begin{proof}
Let $Z_a=Y_a-Y_{d+1}$, $a=1,\ldots,d$. Then $Z$ is Brownian motion
with positive-definite covariance $I+\mathbf1\mathbf1^{\mathsf T}$,
up to its initial translation. It is an invertible linear transform
of standard $d$-dimensional Brownian motion. The map from $Z$ to the
simplex interior is a homeomorphism with inverse
$Z_a=\log(x_a/x_{d+1})$.

We use the classical Brownian recurrence/transience facts
\cite[Section 3.2]{mp}. Their relevant hitting argument is as follows.
For a standard Brownian motion started at radius $r$ between radii
$a$ and $R$, the probability of hitting the inner sphere first is
$(R-r)/(R-a)$ in dimension one, $\log(R/r)/\log(R/a)$ in dimension
two, and $(r^{2-d}-R^{2-d})/(a^{2-d}-R^{2-d})$ in dimension at least
three. These follow by optional stopping of the bounded harmonic
radial functions on the annulus. Letting $R\to\infty$ shows that
every ball is hit almost surely in dimensions one and two. Apply
the Markov property after each deterministic integer time, and use
a countable base of balls, to obtain visits to every open ball at
arbitrarily large times. This property survives an invertible linear
map. The homeomorphism therefore makes every interior strategy a
cluster point; compactness and closedness of the omega set give
$\Omega=X$.

In dimension at least three, the probability of ever returning to
radius $a$ from radius $2^k$ is $(a/2^k)^{d-2}$. The first hitting
times of the spheres $2^k$ are finite almost surely, since one
Brownian coordinate is unbounded. The return probabilities are summable;
the strong Markov property and Borel--Cantelli imply eventual escape
from the ball of radius $a$. Taking integer $a$ proves escape from
every compact set, also after an invertible linear map. Every set of
strategies a positive distance from the simplex boundary has bounded
log ratios. Consequently the strategy path eventually leaves each
such set, so all its cluster points are on the boundary. Compactness
of $X$ supplies at least one cluster point.
\end{proof}

\section*{The unresolved definition and extension}
For $d=1,2$ the whole strategy space is the omega set even though
every proper deterministic stable set is boundary-contained. This
illustrates why the proper-set condition and the precise meaning of
``stable limit set'' matter. For $d\ge3$, finite-time irreducibility
coexists with eventual approach to the boundary. We have not proved
stability of an arbitrary random omega set, nor specified a universally
appropriate random-set stability definition. State-dependent noise,
nonzero strategic drift in the random-set calculation, and the broader
source agenda remain unresolved here. The deterministic benchmark is
explicitly treated as a benchmark already amenable to irreducibility,
rather than advertised as the source conjecture itself.
\begin{thebibliography}{9}
\bibitem{clm} P.-L. Cauvin, D. Legacci, and P. Mertikopoulos (2025),
\emph{The impact of uncertainty on regularized learning in games},
arXiv:2506.13286v1, Sections 2.3,3.1,6.
\url{https://arxiv.org/html/2506.13286v1}.
\bibitem{mp} P. M\"orters and Y. Peres (2010), \emph{Brownian Motion},
Cambridge University Press, Sections 1.4 and 3.2. Author-hosted text:
\url{https://www.mi.uni-koeln.de/~moerters/book/book.pdf}.
\end{thebibliography}

\end{document}
